\scanpage{12}
Noyau $K$ sur le site $X$. Question : $K$ provient-il d'un faisceau en groupes
\[
\text{i.e. }c(K)\in H^2(X,\unsure{I|Z}(K))\text{ est-il dans }\deleted{\unsure{I(K)}}H^2(X,I(K))'\text{ ???}
\]

\deleted{Supposons que l'on puisse trouver un sous-}

À un noyau $K$, est associé le \uline{faisceau} $\unsure{\sigma|\mathscr{T}}(K)$ des \uline{classes} des \uline{sous-groupes de $K$}. \deleted{Supposons \unsure{que nous disposions} d'une section de $\sigma(K)$.}
\inserted{\unsure{On a des opérations} $\sigma(K)\longrightarrow\sigma(K)$, \unsure{telle|t.q.} $\unsure{\varphi|\phi}\longmapsto N(\varphi)$ (\unsure{normalisateur|normaliseur|normalisant}).}
\deleted{Alors on peut \unsure{considérer} le groupe \unsure{normalisateur} dans $\sigma(K)$.}

Chaque fois qu'on a un hom de noyaux $K_0-K$ \deleted{qui est ``injectif''}, il définit une section de $\sigma(K)$ [noyau image].

\deleted{D'ailleurs} Mais une section de $\sigma(K)$ ne définit en général pas un sous-noyau de $K$, sauf celles qui sont stables par l'application $\unsure{\operatorname{Nor}|\operatorname{Na}}$ [``Normalisateur'']. Plus \unsure{généralement|précisément}, si $\varphi$ est une section de $\sigma(K)$, alors $\operatorname{Nor}(\varphi)$ provient d'un sous-noyau de $K$, soit $\underline{\operatorname{Nor}}(\varphi)$, contenant $\unsure{Z|\mathfrak{z}}(K)=\underline{\operatorname{Nor}}(e)$.

\deleted{Supposons que} De plus, on trouve un élément canonique dans $H^2(X,\underline{\operatorname{Nor}}(\varphi))$ dont l'image dans $H^2(X,I(K))$ existe et est $c(K)$ : \deleted{\unsure{il est défini}} c'est la classe de la \unsure{catégorie-loc}
% archive sheet number: 11

